% GENERATED FILE. Edit canonical lessons, then run npm run generate:books.
% canonical-source-hash: fnv1a64:9c66e0106fa83ad0
% canonical-lessons: LA-C133-aranea, LA-C133-rana, LA-C133-lepus, LA-C133-vulpes, LA-C133-taurus

\chapter{Spider, Frog, Hare, Fox, Bull}
\label{ch:la-spider-frog-hare-fox-bull}
\hlchaptermodality{133}

\begin{chapteropening}
\textbf{By the end of this chapter:} \emph{I can say a spider, a frog, a hare, a fox and a bull.}

Complete the last lesson of chapter 133: I can say a spider, a frog, a hare, a fox and a bull.
\end{chapteropening}

\section[arānea, arāneae]{arānea, arāneae — a spider}

\label{lesson:LA-C133-aranea}



\begin{quote}
Before the new one: say the Latin for an ant, then the Latin for a butterfly.
\end{quote}

\subsection*{You'll want to know: arānea, arāneae}
\textbf{arānea, arāneae} — \textquotedblleft{}a spider\textquotedblright{}.

\begin{cousinweb}
\textbf{arānea} is closely related to Greek \emph{arakhnē}, \textquotedblleft{}spider\textquotedblright{} (Lewis \& Short); whether Latin borrowed it or both took it from a third language is uncertain.
\end{cousinweb}

\begin{grammarlens}[title={What you've built}]
One more word for this chapter.
\end{grammarlens}

\subsection*{Your turn}
\begin{itemize}
\raggedright
  \item \emph{Say it:} \emph{arānea}
  \item \emph{Say it:} \emph{arānea}, once more
  \item \emph{Say it:} say \emph{arānea}
  \item \emph{Recall:} say the Latin for an ant, then the Latin for a butterfly, then say \emph{arānea} again
\end{itemize}

\begin{tcolorbox}[breakable,colback=teal!4,colframe=teal!35!black,title={Before you move on}]
What does arānea, arāneae mean? (\textquotedblleft{}A spider\textquotedblright{}.) Say it once more.
\end{tcolorbox}
\section[rāna, rānae]{rāna, rānae — a frog}

\label{lesson:LA-C133-rana}



\begin{quote}
Before the new one: say the Latin for a butterfly, then the Latin for a spider.
\end{quote}

\subsection*{You'll want to know: rāna, rānae}
\textbf{rāna, rānae} — \textquotedblleft{}a frog\textquotedblright{}.

\begin{cousinweb}
\textbf{rāna} has no sure origin; it is often thought to imitate the frog's croak.
\end{cousinweb}

\begin{grammarlens}[title={What you've built}]
One more word for this chapter.
\end{grammarlens}

\subsection*{Your turn}
\begin{itemize}
\raggedright
  \item \emph{Say it:} \emph{rāna}
  \item \emph{Say it:} \emph{rāna}, once more
  \item \emph{Say it:} say \emph{rāna}
  \item \emph{Recall:} say the Latin for a butterfly, then the Latin for a spider, then say \emph{rāna} again
\end{itemize}

\begin{tcolorbox}[breakable,colback=teal!4,colframe=teal!35!black,title={Before you move on}]
What does rāna, rānae mean? (\textquotedblleft{}A frog\textquotedblright{}.) Say it once more.
\end{tcolorbox}
\section[lepus, leporis]{lepus, leporis — a hare}

\label{lesson:LA-C133-lepus}



\begin{quote}
Before the new one: say the Latin for a spider, then the Latin for a frog.
\end{quote}

\subsection*{You'll want to know: lepus, leporis}
\textbf{lepus, leporis} — \textquotedblleft{}a hare\textquotedblright{}.

\begin{cousinweb}
\textbf{lepus} has no sure origin; Varro already compared a Sicilian Greek \emph{leporis}, \textquotedblleft{}hare\textquotedblright{} (Lewis \& Short), and it may be a loan from a language of the Mediterranean.
\end{cousinweb}

\begin{grammarlens}[title={What you've built}]
One more word for this chapter.
\end{grammarlens}

\subsection*{Your turn}
\begin{itemize}
\raggedright
  \item \emph{Say it:} \emph{lepus}
  \item \emph{Say it:} \emph{lepus}, once more
  \item \emph{Say it:} say \emph{lepus}
  \item \emph{Recall:} say the Latin for a spider, then the Latin for a frog, then say \emph{lepus} again
\end{itemize}

\begin{tcolorbox}[breakable,colback=teal!4,colframe=teal!35!black,title={Before you move on}]
What does lepus, leporis mean? (\textquotedblleft{}A hare\textquotedblright{}.) Say it once more.
\end{tcolorbox}
\section[vulpēs, vulpis]{vulpēs, vulpis — a fox}

\label{lesson:LA-C133-vulpes}



\begin{quote}
Before the new one: say the Latin for a frog, then the Latin for a hare.
\end{quote}

\subsection*{You'll want to know: vulpēs, vulpis}
\textbf{vulpēs, vulpis} — \textquotedblleft{}a fox\textquotedblright{}.

\begin{cousinweb}
\textbf{vulpēs} is inherited from Proto-Italic \emph{*wolpi-}, \textquotedblleft{}fox\textquotedblright{} (EtymDB 2.1); Greek \emph{alōpēx}, \textquotedblleft{}fox\textquotedblright{}, is often compared (Lewis \& Short), but the details are uncertain.
\end{cousinweb}

\begin{grammarlens}[title={What you've built}]
One more word for this chapter.
\end{grammarlens}

\subsection*{Your turn}
\begin{itemize}
\raggedright
  \item \emph{Say it:} \emph{vulpēs}
  \item \emph{Say it:} \emph{vulpēs}, once more
  \item \emph{Say it:} say \emph{vulpēs}
  \item \emph{Recall:} say the Latin for a frog, then the Latin for a hare, then say \emph{vulpēs} again
\end{itemize}

\begin{tcolorbox}[breakable,colback=teal!4,colframe=teal!35!black,title={Before you move on}]
What does vulpēs, vulpis mean? (\textquotedblleft{}A fox\textquotedblright{}.) Say it once more.
\end{tcolorbox}
\section[taurus, taurī]{taurus, taurī — a bull}

\label{lesson:LA-C133-taurus}



\begin{quote}
Before the new one: say the Latin for a hare, then the Latin for a fox.
\end{quote}

\subsection*{You'll want to know: taurus, taurī}
\textbf{taurus, taurī} — \textquotedblleft{}a bull\textquotedblright{}.

\begin{cousinweb}
\textbf{taurus} is usually traced to Proto-Indo-European \emph{*tauros}, \textquotedblleft{}bull\textquotedblright{} (EtymDB 2.1); Greek \emph{tauros}, \textquotedblleft{}bull\textquotedblright{}, is a cousin (Lewis \& Short).
\end{cousinweb}

\begin{grammarlens}[title={What you've built}]
One more word for this chapter.
\end{grammarlens}

\subsection*{Your turn}
\begin{itemize}
\raggedright
  \item \emph{Say it:} \emph{taurus}
  \item \emph{Say it:} \emph{taurus}, once more
  \item \emph{Say it:} say \emph{taurus}
  \item \emph{Recall:} say the Latin for a hare, then the Latin for a fox, then say \emph{taurus} again
\end{itemize}

\begin{tcolorbox}[breakable,colback=teal!4,colframe=teal!35!black,title={Before you move on}]
What does taurus, taurī mean? (\textquotedblleft{}A bull\textquotedblright{}.) Say it once more.
\end{tcolorbox}
